% ECONOMICS_PROBLEM: {"id": "H24-325", "title": "Wealth-tax behavioral responses", "jel_code": "H24", "source_id": "OP-0472"}
% !TeX program = xelatex
\documentclass[11pt,a4paper]{article}
\usepackage[margin=23mm]{geometry}
\usepackage{fontspec}
\setmainfont{Latin Modern Roman}
\usepackage{amsmath,amssymb,mathtools}
\usepackage{xcolor,enumitem,booktabs,longtable,array,xurl,hyperref}
\definecolor{muted}{HTML}{53636C}
\hypersetup{colorlinks=true,urlcolor=blue,linkcolor=blue}
\setlength{\parindent}{0pt}
\setlength{\parskip}{5pt}
\newcommand{\R}{\mathbb R}
\newcommand{\E}{\mathbb E}
\newcommand{\N}{\mathbb N}
\newcommand{\ind}{\mathbf 1}
\newcommand{\Var}{\operatorname{Var}}
\newcommand{\Cov}{\operatorname{Cov}}
\newcommand{\supp}{\operatorname{supp}}
\DeclareMathOperator*{\argmin}{arg\,min}
\DeclareMathOperator*{\argmax}{arg\,max}
\newcommand{\entrymeta}[3]{{\sffamily\small\color{muted}\textbf{\texttt{#1}}\quad|\quad #2\par #3\par}}
\newcommand{\Problemref}[1]{\texttt{#1}}
\newcommand{\JELref}[2]{\texttt{#2} (JEL \texttt{#1})}
\begin{document}
\section*{Wealth-tax behavioral responses}
\textbf{Problem ID:} \texttt{H24-325}\quad \textbf{JEL:} \texttt{H24}\par
% BEGIN ECONOMICS STATEMENT
\label{problem:OP-0472}
{\sffamily\small\color{muted}Original subject group: Taxation and social insurance\par}
\label{definitions:problem:30007754}
\textbf{Audit classification:} Published research agenda. The published full text and metadata are verified. Sections 3.5–3.7 document savings, labor and migration evidence gaps; they do not formulate this complete counterfactual vector. Concealment bounds now require explicit restrictions.
\subsection*{Definitions and precise research target}
Fix an initially resident population, implementation date, horizon \(H<\infty\) and constant-price currency. A regime \(p=(\tau,b,v,e)\) specifies rates, exemptions, asset valuations, enforcement, anticipation information and treatment of emigrants; \(p_0\) is a complete baseline. Let \(Z_i^H(p)=(S_i^H(p),I_i^H(p),M_i^H(p),A_i^H(p),R_i^H(p))\) be integrable potential outcomes for each initial resident: cumulative real saving, business investment, the indicator of emigration by \(H\), concealed net assets at \(H\), and discounted public receipts attributable to that person. State which transfers and capital gains enter saving; report investment separately because it can overlap with saving. Fix whether \(p\) changes one jurisdiction or a complete vector of jurisdictions, and fix an equilibrium selection if market responses are multiple. The target is
\[
 \Delta_H(p,p_0)=\mathbb E[Z_i^H(p)]-\mathbb E[Z_i^H(p_0)].
\]
It is a difference of mean regime outcomes and does not require identifying a joint cross-world distribution. Residents who leave remain in the target population; residence-selected samples estimate a different quantity.

Let \(P_O\) be the observable law of linked ownership, residence and tax records. A prespecified class \(\mathcal M\) includes treatment assignment, equilibrium spillovers, anticipation, private-asset measurement and concealment models. Define the identified set \(\mathcal D(P_O)=\{\Delta_H(m):m\in\mathcal M,\ P_O(m)=P_O\}\). Determine when it is a singleton, or derive sharp coordinate and joint bounds. Calling bounds sharp means that every retained target vector is attained or arbitrarily approached by compatible models. With no restriction on concealed assets, their mean need not have finite bounds; either impose an externally justified asset envelope or report unboundedness. Separate real accumulation, asset-price changes, reporting and relocation. A taxable-wealth elasticity does not identify this vector. The joint target is editorial; the cited review provides relevant findings and documented evidence gaps.
\subsection*{Variants and assumption changes}
\begin{enumerate}
\item \textbf{One-off versus recurring tax (editorial).} Compare a one-time liability fixed by preannouncement wealth with annual taxes on subsequently chosen wealth. Saving and relocation incentives differ even at equal initial liabilities.
\item \textbf{Local versus coordinated intervention (editorial).} Compare changing one jurisdiction with changing all competing jurisdictions. Retain the initial-resident population in both; cross-border substitution cannot be held fixed in the second case without an extra assumption.
\item \textbf{Validated versus unrestricted assets (editorial).} Representative verification audits and bounded valuation errors may shrink \(\mathcal D(P_O)\); ordinary self-reports alone need not bound concealment.
\end{enumerate}
\subsection*{References and source audit}
Advani--Tarrant (2021), published version in LSE Research Online, sections 3.5--3.7 (printed pp. 528--532): relevant evidence gaps verified. The IFS page independently confirms author, year, pages and DOI. The joint identified-set specification is editorial.
\begin{enumerate}\raggedright
\item\hypertarget{definitions:source:30007754:1}{} Arun Advani; Hannah Tarrant. 2021. \emph{Behavioural responses to a wealth tax}. Fiscal Studies 42(3–4), 509–537; IFS author abstract, paragraphs 1–2.
\url{https://ifs.org.uk/journals/behavioural-responses-wealth-tax}.
DOI: \url{https://doi.org/10.1111/1475-5890.12283}.
\item\hypertarget{definitions:source:30007754:2}{} Arun Advani; Hannah Tarrant. 2021. \emph{Behavioural responses to a wealth tax}. Published version; sections3.5–3.7, printed pp.528–532.
\url{https://researchonline.lse.ac.uk/id/eprint/112695/1/1475_5890.12283.pdf}.
DOI: \url{https://doi.org/10.1111/1475-5890.12283}.
\end{enumerate}

\subsection*{Common conventions: Source mathematical conventions}

These conventions apply to each entry unless explicitly replaced.

\textbf{Models and identification.} A primitive model \(m\) specifies all probability laws, feasible choices, information, equilibrium and measurement rules used by its target. The admissible class \(\mathcal M\) is fixed independently of outcomes. If \(P_O(m)\) is its observable probability law and \(\tau(m)\) the target, its identified set at observable law \(P\) is
\[
 \mathcal I_\tau(P)=\{\tau(m):m\in\mathcal M,\ P_O(m)=P\}.
\]
Point identification means that this set is a singleton. An empty set signals incompatibility of data and assumptions. Coordinate bounds need not describe the joint set. Bounds are sharp only if every retained target is attainable or arbitrarily approachable by compatible models. Finite bounds require explicit restrictions on unobserved quantities; unrestricted confounding, hidden wealth, extrapolation or arbitrary equilibrium selection can yield uninformative sets.

\textbf{Causal interventions.} A potential outcome \(Y_i(p)\) is the outcome under a complete policy regime \(p\), including timing, financing, information and market spillovers. The target population and units are fixed before treatment. With interference, use the complete assignment vector or a justified exposure mapping; individual no-interference notation alone cannot identify an economy-wide intervention. A difference of mean potential outcomes does not require the cross-world joint distribution. Conditioning on post-treatment migration, survival, enrollment or employment changes the estimand and needs separate justification.

\textbf{Statistical statements.} An estimator, confidence set, forecast or test specifies a sampling law, model class and loss/coverage criterion. Uniform \((1-\alpha)\)-coverage means \(\inf_{m\in\mathcal M}\Pr_m\{\tau(m)\in C_n\}\geq1-\alpha\), or an explicitly asymptotic version. The whole parameter space trivially has coverage, so informativeness requires a stated width, risk or power objective. A forecasting improvement is relative to a named benchmark, information set and evaluation population.

\textbf{Equilibrium and optimization.} An equilibrium comprises feasible plans, optimality, beliefs where needed, budget constraints and all market/resource clearing. The primitives or a stated benchmark define each correspondence \(\mathcal E(p;m)\). Nonemptiness is assumed or proved, never inferred just from compact policy sets. A selected-equilibrium target uses a declared selection rule; a robust target quantifies over all permitted equilibria. Use suprema or infima unless attainment is established. An \(\varepsilon\)-optimal policy is feasible and within \(\varepsilon\) of the finite optimum value. Report an empty feasible set separately.

\textbf{Welfare and accounting.} Normative utility, interpersonal normalization, social weights, numeraire, discounting and the use of public receipts are primitives. Behavioral utility need not coincide with the welfare criterion. Household welfare includes ownership income and rents once; adding profits again to compensating variation generally double-counts them. A monetary equivalent is defined by an explicit transfer and price convention, with monotonicity and range conditions. Average effects, mechanism contributions, transfers and welfare losses are different targets. Infinite sums require convergence and dynamic feasible plans require terminal or no-Ponzi conditions.

\textbf{Source versions.} A working-paper date, revision date and journal publication year are distinguished. A DOI for a journal version is not automatically the DOI of an earlier manuscript. Locators refer to the stated version and printed pages where specified. A metadata-only check does not verify a claim about a full-text conclusion. Every entry's source subsection narrows the provenance claim accordingly.
% END ECONOMICS STATEMENT
\end{document}
